% Practice list B v3 --- topics of Listas 3-4, QUALIFIER LEVEL (style guide §8: cross-topic,
% parameter analysis, graphs, answer key up to 15 pages).
% \solucoesfalse -> blank list ; \solucoestrue -> answer key + rubric.
% Compile: pdflatex (2x); check: pdffonts (all Type 1). Every number: practice-L3-L4-check.py
\documentclass[11pt,a4paper]{article}
% Shared preamble of the practice lists (practice-L*.tex), adapted from lista-exam.tex.
% Each list sets \solucoesfalse / \solucoestrue and \listatema right after % Shared preamble of the practice lists (practice-L*.tex), adapted from lista-exam.tex.
% Each list sets \solucoesfalse / \solucoestrue and \listatema right after % Shared preamble of the practice lists (practice-L*.tex), adapted from lista-exam.tex.
% Each list sets \solucoesfalse / \solucoestrue and \listatema right after \input{practice-preamble}.
\usepackage[utf8]{inputenc}
\usepackage[T1]{fontenc}
\usepackage{lmodern}   % Type 1 vector fonts (replaces PK bitmaps)
\usepackage{cmap}      % ToUnicode CMap -> copyable text and accents
\usepackage[english]{babel}
\usepackage{amsmath,amssymb,amsfonts}
\usepackage[a4paper,margin=2.2cm]{geometry}
\usepackage{enumitem}
\usepackage{fancyhdr}
\usepackage{comment}
\usepackage{xcolor}
\usepackage{microtype}
\usepackage{booktabs}
\usepackage{tikz}
\usepackage{pgfplots}
\pgfplotsset{compat=1.16}

\newif\ifsolucoes
\newcommand{\disciplina}{Macroeconomics I --- MPE Insper}
\newcommand{\espacoresol}{4.5cm}   % height of the answer space in the blank list
% The sol environment must be decided after the list sets its toggle.
\AtBeginDocument{\ifsolucoes\includecomment{sol}\else\excludecomment{sol}\fi}

\pagestyle{fancy}\fancyhf{}
\lhead{\footnotesize\disciplina}
\rhead{\footnotesize Practice list\ifsolucoes{} --- Solutions and rubric\fi}
\cfoot{\footnotesize\thepage}
\renewcommand{\headrulewidth}{0.3pt}
\setlength{\headheight}{14pt}
\setlist{topsep=2pt,itemsep=2pt}
\newcounter{questao}

\newcommand{\bloco}[3]{\bigskip\noindent
  {\large\bfseries\color{blue!55!black}Block #1 --- #2}\par\noindent\hfill\textit{(#3)}\par
  \noindent\rule{\linewidth}{0.3pt}\smallskip}
\newcommand{\questao}[1]{\refstepcounter{questao}\medskip\noindent
  \textbf{Question \thequestao.}~#1\par}
\newenvironment{alternativas}{\begin{enumerate}[label=\Alph*.,leftmargin=2.2em,topsep=1pt]}{\end{enumerate}}
\newcommand{\gabarito}[1]{\ifsolucoes\par\noindent\textbf{Answer:}~#1\fi}
\newcommand{\resolespaco}{\ifsolucoes\else\par\textit{Solution:}\vspace{\espacoresol}\par\fi}
\newcommand{\passo}[2]{\par\smallskip\noindent\textbf{Step --- #1.}~#2}
\newcommand{\intuicao}[1]{\par\smallskip\noindent\textit{Economic intuition.}~#1}
\newcommand{\refbib}[1]{\par\smallskip\noindent{\footnotesize\color{black!60}\textbf{Ref:} #1}}
\newcommand{\rubrica}[1]{\par\smallskip\noindent\textbf{Grading rubric.}\begin{itemize}[leftmargin=1.4em,topsep=1pt]#1\end{itemize}}

\newcommand{\cabecalho}{%
\begin{center}
  {\Large\bfseries \disciplina}\\[2pt]
  {\large Practice list: \listatema}\ifsolucoes\\[2pt]\textit{Solutions and grading rubric}\fi
\end{center}
\vspace{2pt}
{\footnotesize\noindent\textbf{Instructions.} Multiple choice: mark the single correct option (no
penalty for wrong answers). \emph{True/False}: every 2 wrong items cancel 1 right one
(score $=\max\{0,\text{right}-\text{wrong}/2\}\times$ item value); blank items do not count.
Discursive questions: show every step. Numbers are illustrative and chosen to give clean answers.\par}}
.
\usepackage[utf8]{inputenc}
\usepackage[T1]{fontenc}
\usepackage{lmodern}   % Type 1 vector fonts (replaces PK bitmaps)
\usepackage{cmap}      % ToUnicode CMap -> copyable text and accents
\usepackage[english]{babel}
\usepackage{amsmath,amssymb,amsfonts}
\usepackage[a4paper,margin=2.2cm]{geometry}
\usepackage{enumitem}
\usepackage{fancyhdr}
\usepackage{comment}
\usepackage{xcolor}
\usepackage{microtype}
\usepackage{booktabs}
\usepackage{tikz}
\usepackage{pgfplots}
\pgfplotsset{compat=1.16}

\newif\ifsolucoes
\newcommand{\disciplina}{Macroeconomics I --- MPE Insper}
\newcommand{\espacoresol}{4.5cm}   % height of the answer space in the blank list
% The sol environment must be decided after the list sets its toggle.
\AtBeginDocument{\ifsolucoes\includecomment{sol}\else\excludecomment{sol}\fi}

\pagestyle{fancy}\fancyhf{}
\lhead{\footnotesize\disciplina}
\rhead{\footnotesize Practice list\ifsolucoes{} --- Solutions and rubric\fi}
\cfoot{\footnotesize\thepage}
\renewcommand{\headrulewidth}{0.3pt}
\setlength{\headheight}{14pt}
\setlist{topsep=2pt,itemsep=2pt}
\newcounter{questao}

\newcommand{\bloco}[3]{\bigskip\noindent
  {\large\bfseries\color{blue!55!black}Block #1 --- #2}\par\noindent\hfill\textit{(#3)}\par
  \noindent\rule{\linewidth}{0.3pt}\smallskip}
\newcommand{\questao}[1]{\refstepcounter{questao}\medskip\noindent
  \textbf{Question \thequestao.}~#1\par}
\newenvironment{alternativas}{\begin{enumerate}[label=\Alph*.,leftmargin=2.2em,topsep=1pt]}{\end{enumerate}}
\newcommand{\gabarito}[1]{\ifsolucoes\par\noindent\textbf{Answer:}~#1\fi}
\newcommand{\resolespaco}{\ifsolucoes\else\par\textit{Solution:}\vspace{\espacoresol}\par\fi}
\newcommand{\passo}[2]{\par\smallskip\noindent\textbf{Step --- #1.}~#2}
\newcommand{\intuicao}[1]{\par\smallskip\noindent\textit{Economic intuition.}~#1}
\newcommand{\refbib}[1]{\par\smallskip\noindent{\footnotesize\color{black!60}\textbf{Ref:} #1}}
\newcommand{\rubrica}[1]{\par\smallskip\noindent\textbf{Grading rubric.}\begin{itemize}[leftmargin=1.4em,topsep=1pt]#1\end{itemize}}

\newcommand{\cabecalho}{%
\begin{center}
  {\Large\bfseries \disciplina}\\[2pt]
  {\large Practice list: \listatema}\ifsolucoes\\[2pt]\textit{Solutions and grading rubric}\fi
\end{center}
\vspace{2pt}
{\footnotesize\noindent\textbf{Instructions.} Multiple choice: mark the single correct option (no
penalty for wrong answers). \emph{True/False}: every 2 wrong items cancel 1 right one
(score $=\max\{0,\text{right}-\text{wrong}/2\}\times$ item value); blank items do not count.
Discursive questions: show every step. Numbers are illustrative and chosen to give clean answers.\par}}
.
\usepackage[utf8]{inputenc}
\usepackage[T1]{fontenc}
\usepackage{lmodern}   % Type 1 vector fonts (replaces PK bitmaps)
\usepackage{cmap}      % ToUnicode CMap -> copyable text and accents
\usepackage[english]{babel}
\usepackage{amsmath,amssymb,amsfonts}
\usepackage[a4paper,margin=2.2cm]{geometry}
\usepackage{enumitem}
\usepackage{fancyhdr}
\usepackage{comment}
\usepackage{xcolor}
\usepackage{microtype}
\usepackage{booktabs}
\usepackage{tikz}
\usepackage{pgfplots}
\pgfplotsset{compat=1.16}

\newif\ifsolucoes
\newcommand{\disciplina}{Macroeconomics I --- MPE Insper}
\newcommand{\espacoresol}{4.5cm}   % height of the answer space in the blank list
% The sol environment must be decided after the list sets its toggle.
\AtBeginDocument{\ifsolucoes\includecomment{sol}\else\excludecomment{sol}\fi}

\pagestyle{fancy}\fancyhf{}
\lhead{\footnotesize\disciplina}
\rhead{\footnotesize Practice list\ifsolucoes{} --- Solutions and rubric\fi}
\cfoot{\footnotesize\thepage}
\renewcommand{\headrulewidth}{0.3pt}
\setlength{\headheight}{14pt}
\setlist{topsep=2pt,itemsep=2pt}
\newcounter{questao}

\newcommand{\bloco}[3]{\bigskip\noindent
  {\large\bfseries\color{blue!55!black}Block #1 --- #2}\par\noindent\hfill\textit{(#3)}\par
  \noindent\rule{\linewidth}{0.3pt}\smallskip}
\newcommand{\questao}[1]{\refstepcounter{questao}\medskip\noindent
  \textbf{Question \thequestao.}~#1\par}
\newenvironment{alternativas}{\begin{enumerate}[label=\Alph*.,leftmargin=2.2em,topsep=1pt]}{\end{enumerate}}
\newcommand{\gabarito}[1]{\ifsolucoes\par\noindent\textbf{Answer:}~#1\fi}
\newcommand{\resolespaco}{\ifsolucoes\else\par\textit{Solution:}\vspace{\espacoresol}\par\fi}
\newcommand{\passo}[2]{\par\smallskip\noindent\textbf{Step --- #1.}~#2}
\newcommand{\intuicao}[1]{\par\smallskip\noindent\textit{Economic intuition.}~#1}
\newcommand{\refbib}[1]{\par\smallskip\noindent{\footnotesize\color{black!60}\textbf{Ref:} #1}}
\newcommand{\rubrica}[1]{\par\smallskip\noindent\textbf{Grading rubric.}\begin{itemize}[leftmargin=1.4em,topsep=1pt]#1\end{itemize}}

\newcommand{\cabecalho}{%
\begin{center}
  {\Large\bfseries \disciplina}\\[2pt]
  {\large Practice list: \listatema}\ifsolucoes\\[2pt]\textit{Solutions and grading rubric}\fi
\end{center}
\vspace{2pt}
{\footnotesize\noindent\textbf{Instructions.} Multiple choice: mark the single correct option (no
penalty for wrong answers). \emph{True/False}: every 2 wrong items cancel 1 right one
(score $=\max\{0,\text{right}-\text{wrong}/2\}\times$ item value); blank items do not count.
Discursive questions: show every step. Numbers are illustrative and chosen to give clean answers.\par}}

\solucoesfalse                 % <<<< \solucoestrue for the ANSWER KEY
\newcommand{\listatema}{Consumption, Saving and Labor Markets --- qualifier level (Listas 3--4)}
\renewcommand{\espacoresol}{6.5cm}

\begin{document}
\cabecalho

% =============================================================================
\bloco{1}{Consumption, present bias and interest rates}{Kurlat chs.~6, 9; Romer ch.~8 \,\textbullet\, 25 points}

\questao{(10 points) Objective items.}

\textbf{a)} (3 points) A household lives 40 periods ($t=0,\dots,39$) with CRRA utility, $\sigma=2$, $\beta=0.97$, faces $r=4\%$
and has total wealth $W$ (financial plus human, in present value at $t=0$). The share of $W$ consumed at $t=0$ is:
\begin{alternativas}
\item $3.42\%$
\item $4.55\%$
\item $2.50\%$
\item $3.85\%$
\end{alternativas}
\gabarito{B}
\begin{sol}
\passo{path}{Euler: $c_t=c_0[\beta(1+r)]^{t/\sigma}$; budget $\sum_{t=0}^{39}c_t(1+r)^{-t}=W$ gives $c_0\sum_{t=0}^{39}\gamma^t=W$ with
$\gamma\equiv\beta^{1/\sigma}(1+r)^{1/\sigma-1}=\sqrt{0.97/1.04}=0.96576$.}
\passo{share}{$c_0/W=\frac{1-\gamma}{1-\gamma^{40}}=\frac{0.03424}{0.75182}=4.55\%$. A is the infinite-horizon $1-\gamma$; C is $1/40$; D is $r/(1+r)$.}
\refbib{Kurlat (2020), \S6.3, eqs.\ (6.3.6)--(6.3.9), pp.~117--120.}
\end{sol}

\medskip
\textbf{b)} (3 points) \emph{Consumption $\times$ general equilibrium.} A two-period endowment economy with a
representative household: $y_2/y_1=1.03$, $\beta=0.97$, CRRA with $\sigma=2$. The equilibrium real interest rate is:
\begin{alternativas}
\item $3.09\%$
\item $6.19\%$
\item $9.37\%$
\item $3.00\%$
\end{alternativas}
\gabarito{C}
\begin{sol}
\passo{clearing}{No one can lend in the aggregate: $c_t=y_t$. Euler: $y_1^{-\sigma}=\beta(1+r)y_2^{-\sigma}\Rightarrow1+r=(y_2/y_1)^{\sigma}/\beta=1.0609/0.97$,
$r=9.37\%$. B uses $\sigma=1$; A ignores growth.}
\refbib{Kurlat (2020), \S9.1, pp.~165--168.}
\end{sol}

\medskip
\textbf{c)} (2 points, T or F) \emph{Consumption $\times$ credit markets.} A household strictly inside the region where it
neither borrows nor lends (borrowing rate above the deposit rate) starts saving after a small rise in the deposit
rate.
\gabarito{F}
\begin{sol}
\passo{kink}{At the kink the MRS $u'(y_1)/[\beta u'(y_2)]$ lies strictly between $1+r_s$ and $1+r_b$; a small rise in $r_s$ keeps it
above $1+r_s$, so saving is still unattractive and consumption does not change.}
\refbib{Kurlat (2020), Exercise 6.6 (kinked budget sets).}
\end{sol}

\medskip
\textbf{d)} (2 points, T or F) \emph{Consumption $\times$ fiscal policy.} A household at a binding borrowing limit receives a
lump-sum tax cut of 100 today, financed by a tax increase of $100(1+r)$ next period. If the limit still binds after
the cut, current consumption rises by exactly 100.
\gabarito{T}
\begin{sol}
\passo{why}{With the limit binding, $c_1=y_1-\tau_1+b$: consumption equals current resources, so it moves one for one with
the tax cut; Ricardian equivalence fails because the household cannot borrow against the future tax.}
\refbib{Kurlat (2020), \S6.2 and Exercise 6.5.}
\end{sol}

% -----------------------------------------------------------------------------
\questao{(15 points) Discursive. \emph{Consumption $\times$ behavioral economics $\times$ general equilibrium.} Three periods
$t=0,1,2$; in (a)--(c) the interest rate is zero and the household has wealth $W$ at date 0. The self at date $t$ maximizes
$u(c_t)+\beta\sum_{s>t}u(c_s)$ with $u(c)=c^{1-\sigma}/(1-\sigma)$ ($\ln c$ if $\sigma=1$) and $0<\beta<1$.}

\textbf{a)} (3 points) Find the plan chosen with commitment (self 0 chooses all three consumptions), by a naive self 0
(who believes future selves will follow its plan), and by a sophisticated self 0 (who anticipates self 1's choice).
\resolespaco
\begin{sol}
\passo{self 1}{$\max u(c_1)+\beta u(c_2)$ s.t.\ $c_1+c_2=W_1$: $c_1^{-\sigma}=\beta c_2^{-\sigma}\Rightarrow c_2=\beta^{1/\sigma}c_1$, so $c_1=kW_1$, $c_2=(1-k)W_1$,
$k\equiv\frac1{1+\beta^{1/\sigma}}>\frac12$.}
\passo{commitment}{Self 0 weighs dates 1 and 2 equally: $c_1=c_2=c$, and $u'(c_0)=\beta u'(c)\Rightarrow c_0=\beta^{-1/\sigma}c$; the budget gives
$c=\frac{W}{2+\beta^{-1/\sigma}}$.}
\passo{naive}{Believing $c_1=c_2$, self 0 picks the commitment $c_0$; then self 1 re-optimizes: $c_1=kW_1$, $c_2=(1-k)W_1$.}
\passo{sophisticated}{Self 0 solves $\max_{c_0}u(c_0)+\beta[u(kW_1)+u((1-k)W_1)]$, $W_1=W-c_0$. FOC:
$c_0^{-\sigma}=\beta[k^{1-\sigma}+(1-k)^{1-\sigma}]W_1^{-\sigma}$, so $c_0=\frac{x_s}{1+x_s}W$ with $x_s\equiv\big(\beta[k^{1-\sigma}+(1-k)^{1-\sigma}]\big)^{-1/\sigma}$ (for $\sigma=1$,
$x_s=1/(2\beta)$).}
\refbib{Romer (2012), \S8.6 and problem 8.15, pp.~397--404; Kurlat (2020), \S6.4.}
\rubrica{
\item Self 1's rule: 0.5 pt.
\item Commitment plan: 1 pt.
\item Naive plan (two steps): 0.5 pt.
\item Sophisticated FOC and $c_0$: 1 pt.
}
\end{sol}

\textbf{b)} (3 points) Determine when the sophisticated self 0 consumes more, less or the same as the naive one,
and explain the mechanism.
\resolespaco
\begin{sol}
\passo{compare}{Naive: $c_0/W_1=x_n=\beta^{-1/\sigma}/2$, so $x_n^{-\sigma}=\beta2^{\sigma}=2\beta\varphi(\tfrac12)$ with $\varphi(z)\equiv z^{1-\sigma}$. Sophisticated:
$x_s^{-\sigma}=\beta[\varphi(k)+\varphi(1-k)]$. Since $z\mapsto z^{-\sigma}$ is decreasing and $c_0=\frac{x}{1+x}W$ rises with $x$:
$c_0^s<c_0^n\iff\varphi(k)+\varphi(1-k)>2\varphi(\tfrac12)$.}
\passo{Jensen}{$\varphi''(z)=\sigma(\sigma-1)z^{-\sigma-1}$: $\varphi$ is strictly convex if $\sigma>1$, linear in logs at $\sigma=1$ (both sides equal),
concave if $\sigma<1$. With $k\ne\frac12$: $\sigma>1\Rightarrow c_0^s<c_0^n$; $\sigma=1\Rightarrow$ equal; $\sigma<1\Rightarrow c_0^s>c_0^n$.}
\passo{mechanism}{The sophisticate knows the wealth it leaves will be split unevenly. With high curvature ($\sigma>1$) an
uneven split raises the marginal value of extra wealth, so it saves more to compensate; with $\sigma<1$ the
misallocation lowers that value and it consumes more now. Under log the two forces cancel exactly.}
\refbib{Romer (2012), problem 8.15(b)(ii).}
\rubrica{
\item Reduction to the comparison of $\varphi$ values: 1.5 pts.
\item Jensen argument with all three cases: 1 pt.
\item Mechanism: 0.5 pt.
}
\end{sol}

\textbf{c)} (3 points) What share of $W$ would self 0 pay for an illiquid account that commits future consumption?
Find a formula, its limits as $\beta\to1$ and as $\sigma\to\infty$, and compute it for $\sigma=2$, $\beta=1/2$.
\resolespaco
\begin{sol}
\passo{homogeneity}{For $\sigma\ne1$ every plan is proportional to $W$, so self 0's value is $V_j(W)=W^{1-\sigma}v_j$ ($j=$ commitment $c$,
sophisticated $s$). The willingness to pay $x$ solves $V_c((1-x)W)=V_s(W)$:
\[x=1-\Big(\frac{v_s}{v_c}\Big)^{\frac1{1-\sigma}};\qquad\sigma=1:\ x=1-\exp\Big(-\frac{V_c-V_s}{1+2\beta}\Big).\]
It does not depend on $W$.}
\passo{limits}{$\beta\to1$: $k\to\frac12$, all plans coincide, $x\to0$. $\sigma\to\infty$: $\beta^{1/\sigma}\to1$, $k\to\frac12$, self 1 also splits evenly, $x\to0$
(the table in (e) shows the decline).}
\passo{numbers}{$\sigma=2$, $\beta=\frac12$, $W=100$, plans $(c_0,c_1,c_2)$:
\[\begin{aligned}&\text{commitment }(41.42,\,29.29,\,29.29),\quad\text{naive }(41.42,\,34.31,\,24.26),\\
&\text{sophisticated }(41.06,\,34.53,\,24.41);\end{aligned}\]
$V_c=-0.05828$, $V_s=-0.05932$, $V_n=-0.05932$; $x=1-V_c/V_s=1.74\%$ (log: $2.90\%$).}
\refbib{Romer (2012), problem 8.15; Kurlat (2020), \S6.4 (commitment, defaults).}
\rubrica{
\item Homogeneity and the formula: 1.5 pts.
\item Limits: 0.5 pt.
\item Numbers: 1 pt.
}
\end{sol}

\textbf{d)} (3 points) Now every household receives $y$ at each date and households are identical and sophisticated;
they can borrow and lend at $r_{01}$ (between dates 0 and 1) and $r_{12}$. Find the equilibrium interest rates, show
whether they depend on $\sigma$, and draw them against $\beta\in[0.4,1]$.
\resolespaco
\begin{sol}
\passo{clearing}{Identical households cannot lend to each other: $c_t=y$ at every date.}
\passo{$r_{12}$}{Self 1's Euler at $c_1=c_2=y$: $y^{-\sigma}=\beta(1+r_{12})y^{-\sigma}\Rightarrow1+r_{12}=1/\beta$.}
\passo{self 1's rule at these prices}{With $\beta(1+r_{12})=1$, self 1 sets $c_2=c_1$ and $c_1=W_1/(1+\beta)$, where $W_1=y+y/(1+r_{12})=y(1+\beta)$
in equilibrium: $c_1=y$. Self 0's value of $W_1$ is $2u(W_1/(1+\beta))$, with marginal value $\frac{2}{1+\beta}u'(y)$.}
\passo{$r_{01}$}{Self 0's Euler: $u'(y)=(1+r_{01})\beta\frac{2}{1+\beta}u'(y)\Rightarrow1+r_{01}=\frac{1+\beta}{2\beta}$. Neither rate depends on $\sigma$ (equal endowments). At
$\beta=\frac12$: $r_{01}=50\%$, $r_{12}=100\%$; with $\beta=1$ both are zero.}
\begin{center}
\begin{tikzpicture}
\begin{axis}[width=9.5cm,height=5cm,xmin=0.4,xmax=1,ymin=0,ymax=1.55,xlabel={$\beta$},ylabel={rate},
  legend style={font=\footnotesize},tick label style={font=\footnotesize},label style={font=\small}]
\addplot[thick,blue,domain=0.4:1,samples=80]{(1+x)/(2*x)-1};
\addplot[thick,red,domain=0.4:1,samples=80]{1/x-1};
\legend{$r_{01}$,$r_{12}$}
\end{axis}
\end{tikzpicture}
\end{center}
\intuicao{Present bias makes everyone want to bring consumption forward; with nobody to borrow from, rates
rise until each self is content with its endowment, and the near-term rate rises most.}
\refbib{Kurlat (2020), \S9.1 (endowment equilibrium), pp.~165--168.}
\rubrica{
\item Market clearing and $r_{12}$: 1 pt.
\item Self 0's marginal value of wealth and $r_{01}$: 1 pt.
\item Graph: axes, two decreasing curves meeting at zero when $\beta=1$, $r_{12}$ above $r_{01}$: 1 pt.
}
\end{sol}

\textbf{e)} (3 points) \emph{Parameter analysis.} How do the value of commitment and the equilibrium rates depend on
$\beta$ and $\sigma$? Prove the signs you can, fill the value of commitment for $\sigma\in\{1,1.5,2,5\}$ and $\beta\in\{0.5,0.7,0.9\}$,
and read the results.
\resolespaco
\begin{sol}
\passo{rates}{$\partial r_{01}/\partial\beta=-\frac1{2\beta^2}<0$, $\partial r_{12}/\partial\beta=-\frac1{\beta^2}<0$, $\partial/\partial\sigma=0$; $r_{12}>r_{01}$ for $\beta<1$.}
\passo{table}{Willingness to pay for commitment (\% of $W$):
\begin{center}
\begin{tabular}{lccc}
\toprule
 & $\beta=0.5$ & $\beta=0.7$ & $\beta=0.9$ \\
\midrule
$\sigma=1$ & 2.90 & 0.92 & 0.09 \\
$\sigma=1.5$ & 2.20 & 0.65 & 0.06 \\
$\sigma=2$ & 1.74 & 0.50 & 0.05 \\
$\sigma=5$ & 0.75 & 0.21 & 0.02 \\
\bottomrule
\end{tabular}
\end{center}
It falls with $\beta$ (less disagreement between selves) and with $\sigma$ (self 1's split approaches the even split).}
\passo{reading}{The demand for commitment is concentrated where present bias is strong and curvature is low:
that is where defaults and illiquid accounts (401(k)s, Brazil's FGTS) matter. In equilibrium the same bias shows
up as high short-term real rates, independent of $\sigma$ when endowments are flat.}
\refbib{Kurlat (2020), \S6.4 and \S9.1; Romer (2012), problem 8.15.}
\rubrica{
\item Signs of the rates: 1 pt.
\item Table: 1 pt.
\item Reading of both dependences: 1 pt.
}
\end{sol}

% =============================================================================
\bloco{2}{Labor supply, taxes, transfers and search}{Kurlat ch.~7; Romer ch.~10 \,\textbullet\, 25 points}

\questao{(10 points) Objective items.}

\textbf{a)} (3 points) \emph{Labor $\times$ consumption smoothing.} Per-period utility $\ln c-\psi\frac{h^{1+1/\varepsilon}}{1+1/\varepsilon}$, $\varepsilon=0.5$; 40 periods with
$\beta(1+r)=1$ and equal wages. The wage rises 10\% in one period only. To first order, hours in that period rise by:
\begin{alternativas}
\item $4.9\%$
\item $5.0\%$
\item $0\%$
\item $3.3\%$
\end{alternativas}
\gabarito{A}
\begin{sol}
\passo{FOCs}{$\psi h_t^{1/\varepsilon}=\lambda w_t$ with $\lambda=1/c$ constant across periods; so $d\ln h_1=\varepsilon(d\ln w_1-d\ln c)$ and
$d\ln h_t=-\varepsilon\,d\ln c$ for $t\ne1$.}
\passo{budget}{$d\ln c=\sum_ts_t(d\ln w_t+d\ln h_t)$ with earnings shares $s_t=1/40$; solving, $d\ln c=s_1d\ln w_1$ and
$d\ln h_1=\varepsilon(1-s_1)d\ln w_1=0.5\cdot0.975\cdot10\%=4.9\%$ (4.6\% for the exact discrete change). B is the pure Frisch
response; C the response to a permanent raise.}
\refbib{Kurlat (2020), \S7.4, pp.~140--141.}
\end{sol}

\medskip
\textbf{b)} (3 points) \emph{Search $\times$ wage setting.} In the steady-state search model with $f=\mu\theta^{\alpha}$, $\alpha=0.5$, free entry
$(y-w)q(\theta)=\kappa(r+s)$, the wage is rigid at $w=0.95$ while productivity $y=1$ moves. The elasticity of tightness $\theta$
with respect to $y$ is:
\begin{alternativas}
\item $3.7$
\item $40$
\item $20$
\item $2.0$
\end{alternativas}
\gabarito{B}
\begin{sol}
\passo{log-differentiate}{$q=\mu\theta^{\alpha-1}$: $\ln(y-w)+(\alpha-1)\ln\theta=$ const, so $\frac{d\ln\theta}{d\ln y}=\frac1{1-\alpha}\frac{y}{y-w}=2\cdot20=40$. A is the
bargained-wage value: rigid wages are one answer to Shimer's puzzle.}
\refbib{Kurlat (2020), \S7.5, p.~146; Romer (2012), problem 10.15.}
\end{sol}

\medskip
\textbf{c)} (2 points, T or F) \emph{Search $\times$ efficiency.} With a general constant-returns matching function, the search
equilibrium is efficient iff the workers' share of the surplus equals the elasticity of matches with respect to
unemployment, evaluated at the equilibrium tightness.
\gabarito{T}
\begin{sol}
\passo{Hosios}{The planner's and the market's conditions coincide iff the firms' share equals the vacancy elasticity
of matches; with constant returns the two elasticities add to one, and both are evaluated where the economy is.}
\refbib{Romer (2012), \S10.7 and problem 10.16.}
\end{sol}

\medskip
\textbf{d)} (2 points, T or F) \emph{Labor $\times$ transfers.} With $u=\ln c+\alpha\ln(1-h)$ and a lump-sum transfer $T>0$ that each household
takes as given, an individual household's hours do not respond to its after-tax wage.
\gabarito{F}
\begin{sol}
\passo{supply}{$h=\frac1{1+\alpha}-\frac{\alpha T}{(1+\alpha)w_n}$ rises with $w_n$ when $T>0$; hours are flat in $w_n$ only when $T=0$ (Question 4(a)).}
\refbib{Kurlat (2020), \S7.2.}
\end{sol}

% -----------------------------------------------------------------------------
\questao{(15 points) Discursive. \emph{Labor supply $\times$ taxes and transfers $\times$ welfare.} Households have
$u=\ln c+\alpha\ln(1-h)$ with one unit of time; a linear technology pays $w=1$ per hour; labor income is taxed at $\tau$ and the
revenue is returned as a lump-sum transfer $T$ that each household takes as given. Baseline $\alpha=0.8$.}

\textbf{a)} (3 points) Derive individual labor supply given $(\tau,T)$. Show what the tax does when $T=0$, and find
equilibrium hours $h(\tau)$ when the budget balances.
\resolespaco
\begin{sol}
\passo{FOC}{$\max\ln[(1-\tau)h+T]+\alpha\ln(1-h)$: $\frac{1-\tau}{c}=\frac{\alpha}{1-h}$, i.e.\ $\alpha c=(1-\tau)(1-h)$ with $c=(1-\tau)h+T$.}
\passo{individual supply}{$\alpha[(1-\tau)h+T]=(1-\tau)(1-h)\Rightarrow h=\frac1{1+\alpha}-\frac{\alpha T}{(1+\alpha)(1-\tau)}$. With $T=0$, $h=\frac1{1+\alpha}$:
the tax's income and substitution effects cancel.}
\passo{balanced budget}{$T=\tau h$ (each household still takes $T$ as given): $c=h$, and the FOC gives
$\frac{1-\tau}{h}=\frac{\alpha}{1-h}\Rightarrow h(\tau)=\frac{1-\tau}{1-\tau+\alpha}$, decreasing in $\tau$; $h(0)=0.556$.}
\refbib{Kurlat (2020), \S7.2--7.3 and Exercise 7.5, pp.~131--140, 148--149.}
\rubrica{
\item FOC: 1 pt.
\item Individual supply and the $T=0$ case: 1 pt.
\item Equilibrium $h(\tau)$ with $T$ substituted after the FOC: 1 pt. Substituting before: max.\ 1.5 pts in the item.
}
\end{sol}

\textbf{b)} (3 points) Find the tax rate that maximizes revenue (the size of the transfer program), prove that revenue
is single-peaked, and draw revenue against $\tau$.
\resolespaco
\begin{sol}
\passo{revenue}{$R(\tau)=\tau h(\tau)=\frac{\tau(1-\tau)}{1+\alpha-\tau}$. The numerator of $R'$ is $N(\tau)=\tau^2-2(1+\alpha)\tau+(1+\alpha)$.}
\passo{single peak}{$N$ is a convex parabola with $N(0)=1+\alpha>0$ and $N(1)=-\alpha<0$: exactly one root in $(0,1)$, the other above 1.
So $R'>0$ before and $R'<0$ after
\[\tau^*=(1+\alpha)-\sqrt{\alpha(1+\alpha)}=1.8-1.2=0.6,\qquad R(\tau^*)=0.6\cdot\tfrac{0.4}{1.2}=0.2.\]}
\begin{center}
\begin{tikzpicture}
\begin{axis}[width=9.5cm,height=5cm,xmin=0,xmax=1,ymin=0,ymax=0.23,xlabel={$\tau$},ylabel={$R(\tau)$},
  tick label style={font=\footnotesize},label style={font=\small}]
\addplot[thick,blue,domain=0:1,samples=100]{x*(1-x)/(1.8-x)};
\addplot[only marks,mark=*] coordinates {(0.6,0.2) (0.4,0.171429) (0.5,0.192308)};
\node[font=\footnotesize,anchor=south] at (axis cs:0.6,0.2) {$\tau^*=0.6$};
\end{axis}
\end{tikzpicture}
\end{center}
\refbib{Kurlat (2020), \S7.3.}
\rubrica{
\item Revenue and the sign of $R'$: 1 pt.
\item Single-peak proof and $\tau^*$: 1 pt.
\item Graph: axes, hump-shaped curve, peak at $(0.6,0.2)$: 1 pt.
}
\end{sol}

\textbf{c)} (3 points) Derive the welfare cost of the program in consumption-equivalent terms (the share of
consumption at $\tau=0$ the household would give up to avoid the program), prove that it is zero to first order at
$\tau=0$ and increasing, and compute it at $\tau=0.4$, $0.5$, $0.6$.
\resolespaco
\begin{sol}
\passo{welfare}{With $c=h$: $U(\tau)=\ln h(\tau)+\alpha\ln(1-h(\tau))$. The cost $x$ solves $U(0)+\ln(1-x)=U(\tau)$:
$x(\tau)=1-e^{U(\tau)-U(0)}$.}
\passo{shape}{$U'(\tau)=h'(\tau)\big[\frac1h-\frac{\alpha}{1-h}\big]$. At $\tau=0$, $h=\frac1{1+\alpha}$ and the bracket is zero: no first-order cost. For
$\tau>0$, $h<\frac1{1+\alpha}$, so $\frac1h>1+\alpha>\frac{\alpha}{1-h}$; with $h'<0$, $U'<0$: the cost rises with $\tau$.}
\passo{numbers}{$\tau=0.4$: $h=0.429$, $x=5.68\%$; $\tau=0.5$: $h=0.385$, $x=10.18\%$; $\tau=0.6$: $h=0.333$, $x=17.01\%$.}
\intuicao{The program moves nothing between households; its whole effect is a wedge on the price of leisure,
and the loss grows with the square of the wedge near zero.}
\refbib{Kurlat (2020), Exercise 7.5 (welfare comparison).}
\rubrica{
\item Welfare measure: 1 pt.
\item First-order zero and monotonicity proved: 1 pt.
\item Numbers: 1 pt.
}
\end{sol}

\textbf{d)} (3 points) At $\tau=0.4$, compute the Marshallian and Hicksian elasticities of individual hours with respect to the
after-tax wage, the balanced-budget elasticity $d\ln h/d\ln(1-\tau)$, and the Frisch elasticity. Which one does Prescott's
experiment use?
\resolespaco
\begin{sol}
\passo{Marshallian}{From (a), $\partial h/\partial w_n=\frac{\alpha T}{(1+\alpha)w_n^2}$, so $e_M=\frac{\alpha T}{(1+\alpha)w_nh}=\frac{\alpha\tau}{(1+\alpha)(1-\tau)}=0.296$ (using $T=\tau h$).}
\passo{Hicksian}{Slutsky: $\partial h^c/\partial w_n=\partial h/\partial w_n-h\,\partial h/\partial T=\frac{\alpha c}{(1+\alpha)w_n^2}$, so $e_H=\frac{\alpha}{(1+\alpha)(1-\tau)}=0.741$.}
\passo{balanced budget and Frisch}{$\frac{d\ln h}{d\ln(1-\tau)}=1-h=0.571$. Frisch (marginal utility of wealth fixed): $\frac{\alpha}{1-h}=\lambda w_n$ gives
$\frac{1-h}{h}=1.333$.}
\passo{which one}{Prescott compares economies whose revenue is rebated, so the relevant response is the
balanced-budget one: larger than the Marshallian (the transfer offsets part of the income effect) and smaller
than the Hicksian.}
\refbib{Kurlat (2020), \S7.2--7.3 and footnote 8, p.~149.}
\rubrica{
\item Each elasticity derived: 0.5 pt each (2 pts).
\item Which one and why: 1 pt.
}
\end{sol}

\textbf{e)} (3 points) \emph{Parameter analysis.} How do $\tau^*$, hours at $\tau=0.5$ and the welfare cost at $\tau=0.5$ depend on $\alpha$?
Prove the signs you can, give the limits of $\tau^*$, fill the table for $\alpha\in\{0.4,0.8,1.54,3\}$ (1.54 is the value in Kurlat's
Exercise 7.5) and read it.
\resolespaco
\begin{sol}
\passo{signs}{$\frac{d\tau^*}{d\alpha}=1-\frac{1+2\alpha}{2\sqrt{\alpha(1+\alpha)}}<0$ because $(1+2\alpha)^2=4\alpha(1+\alpha)+1$. Hours $\frac{1-\tau}{1-\tau+\alpha}$ fall with $\alpha$.}
\passo{limits}{$\alpha\to0$: $\tau^*\to1$ (leisure worthless, no distortion until everything is taxed). $\alpha\to\infty$: $\tau^*\to\frac12$, from
$\sqrt{\alpha^2+\alpha}=\alpha+\frac12-\frac1{8\alpha}+\dots$}
\passo{table}{
\begin{center}
\begin{tabular}{lccc}
\toprule
 & $\tau^*$ & $h(0.5)$ & cost at $\tau=0.5$ \\
\midrule
$\alpha=0.4$ & 0.652 & 0.556 & 7.19\% \\
$\alpha=0.8$ & 0.600 & 0.385 & 10.18\% \\
$\alpha=1.54$ & 0.562 & 0.245 & 12.75\% \\
$\alpha=3$ & 0.536 & 0.143 & 14.70\% \\
\bottomrule
\end{tabular}
\end{center}}
\passo{reading}{A higher weight on leisure makes hours more responsive at the margin: the Laffer peak arrives earlier
and the same 50\% wedge costs more. Prescott's account of Europe needs a large $\alpha$, i.e.\ an elasticity at the top
of the micro estimates (Kurlat \S7.3).}
\refbib{Kurlat (2020), \S7.3 and Exercise 7.5.}
\rubrica{
\item Sign of $d\tau^*/d\alpha$ proved and both limits: 1.5 pts.
\item Table: 1 pt.
\item Reading: 0.5 pt.
}
\end{sol}

\end{document}
